\domainevide{Les cotes de ce thème supposent les correspondances algébriques et une catégorie de motifs, que mathlib ne définit pas.}{12, 162-6}

\begin{encadre}
\begin{theoreme*}[Grothendieck, pages 115 à 119, poids et niveau d'une immersion\piste{12-immersion-level-bound}]\nonverifie
Soient $f : X \to Y$ une immersion de source régulière et $n$ la dimension de l'adhérence de $f(X)$. Les constituants de $R^i f_*\mathbb Q_\ell$ sont de poids $\leq 2\inf(i, n)$ et de niveau $\leq \inf(i, n-1)$. Le dossier tient l'énoncé sous réserve de la résolution, de la pureté et des conjectures de Weil.
\end{theoreme*}
\end{encadre}
{\small Les feuillets en esquissent une démonstration incomplète (pages 115, 117 et 119). Un constituant de poids $w$ est de niveau $\leq \nu$ quand ses valeurs propres de Frobenius sont divisibles par $q^{(w-\nu)/2}$, ou, côté Hodge, quand ses types $(p, q)$ vérifient $p, q \geq (w - \nu)/2$ ; la borne $\inf(i, n-1)$ est inférieure d'une unité au niveau $2n - i$ que permet $H^i$ d'une variété lisse de dimension~$n$ quand $i \geq n$. L'édition vérifie que $n - 1$ est atteint : pour $Y$ le cône affine sur une courbe $C$ de genre $\geq 1$ et $X = Y$ privé du sommet ($n = 2$), la fibre de $R^2f_*$ au sommet est $H^1(C)(-1)$, de poids~$3$ et de niveau~$1$. Les tableaux qui mènent à l'encadré sont écrits vite et surchargés, et l'énoncé repose sur l'encadré seul ; le passage des termes $E_2$ à l'aboutissement (proposition~4.2~a) de la page~112) n'est pas écrit. Les structures de Hodge et les périodes sont l'objet de la partie \q{Théorie de Hodge} (cotes~13 et~18). Mathlib n'a pas les poids des faisceaux $\ell$-adiques. On se reportera aux notes de Grothendieck (\lot{12}{6}{115}).\par}
